\chapter{Execution: order types and market impact}\label{ch:exec}

\noindent Alpha decides what to trade; execution decides what it costs to trade it. The
two are the same trade: a strategy that needs to lift 10\% of ADV in an hour is a
different product than its backtest said. This chapter is the mechanics: order types,
the square-root impact law, and the scheduling problem.

\section{Order types}\label{sec:ordtypes}

Three primitives cover the menu: \emph{market} (fill now, pay spread plus impact),
\emph{limit} (fill maybe, pay nothing and sometimes earn a rebate), \emph{scheduled}
(slices over time, paying a mix). Everything else---iceberg, pegged, stop---is a wrapper
around these three trade-offs.

\begin{figure}[htbp]\centering
\begin{tikzpicture}[font=\scriptsize,every node/.style={align=center}]
  \node[draw=ink,fill=goodbg,rounded corners=1pt] (q) at (0,2.2) {order to route};
  \node[draw=slate,fill=white,rounded corners=1pt] (m) at (-3.0,0.8) {market:\\ now,\\ pay spread+impact};
  \node[draw=slate,fill=white,rounded corners=1pt] (l) at (0,0.8) {limit:\\ maybe,\\ free+rebate};
  \node[draw=slate,fill=white,rounded corners=1pt] (s) at (3.0,0.8) {scheduled:\\ TWAP/VWAP,\\ trade speed vs risk};
  \draw[->,ink] (q.south) -- (m.north) node[midway,left]{urgent};
  \draw[->,ink] (q.south) -- (l.north) node[midway,right]{patient};
  \draw[->,ink] (q.south) -- (s.north) node[midway,right]{large};
\end{tikzpicture}

\caption{The routing triage: urgent $\to$ market, patient $\to$ limit, large $\to$ schedule.}
\label{fig:orders}
\end{figure}

\begin{example}\label{ex:mktlim}
\textbf{The trade-off in one line.} A market buy at spread 2bp crosses the spread now;
a limit buy at the bid waits and may never fill.
\begin{solution}
The market order pays $\frac{1}{2}$-spread plus impact for certainty; the limit order
earns the half-spread for patience, accepting adverse selection---you get filled
precisely when the market wants to sell to you, i.e.\ before it moves against you.
Neither is free; they pay in different currencies (certainty vs.\ waiting+selection).
\end{solution}
\end{example}

\section{The square-root law}\label{sec:sqrtlaw}

Empirical market impact is concave: moving $q$ shares costs approximately
$c(q)\approx\sigma\sqrt{q/V}$, where $V$ is daily volume. \texttt{impact\_model.py}:
\begin{lstlisting}[numbers=none]
q=   1000 (0.1% of ADV)  impact    6.3bp
q=  10000 (1.0% of ADV)  impact   20.0bp
q= 100000 (10.0% of ADV)  impact   63.2bp
q= 500000 (50.0% of ADV)  impact  141.4bp
\end{lstlisting}

\begin{example}\label{ex:sqrt}
\textbf{Why the square root matters.} Doubling size multiplies cost by $\sqrt2\approx1.41$,
not 2.
\begin{solution}
Concavity is the mathematical reason big orders get sliced: the marginal cost is highest
for the largest prints, so spreading the same notional over more, smaller prints lowers
total impact. It is also why the impact estimate is the first question asked of any
strategy's capacity---a 20bp edge on a 63bp print is not a product.
\end{solution}
\end{example}

\begin{figure}[htbp]\centering
\begin{tikzpicture}
\begin{axis}[width=0.86\linewidth,height=0.34\linewidth,
    xlabel={order size (\% ADV)},ylabel={impact (bp)},
    xmin=0,xmax=60,ymin=0,ymax=150,
    tick style={draw=none},axis line style={draw=slate},font=\scriptsize]
  \addplot[accent,very thick,domain=0.1:60,samples=80] {20*sqrt(x)};
  \addplot[red,dashed,thick] coordinates {(10,0)(10,63.2)};
  \node[font=\scriptsize,text=red] at (axis cs:16,55) {10\% ADV $\to$ 63bp};
\end{axis}
\end{tikzpicture}

\caption{Square-root impact: cost per unit grows slower than size, but the cost-to-edge
ratio grows without bound.}
\label{fig:impact}
\end{figure}

\section{Slicing: TWAP vs VWAP}\label{sec:slicing}

Given $Q$ to execute over $T$ slices, TWAP spreads evenly; VWAP spreads proportional to
expected volume. \texttt{twap\_vwap.py} on a U-shaped volume curve:
\begin{lstlisting}[numbers=none]
slices 48  volume U-shaped (open/close heavy)
IS TWAP:  -78.10bp   IS VWAP:  -72.72bp
vwap wins on U-volume: schedules weight to where liquidity is
\end{lstlisting}
(The shortfall is negative here because the random walk drifted in the buyer's favour---the
\emph{comparison} is the point, not the sign.)

\begin{example}\label{ex:vwap}
\textbf{Why VWAP beats TWAP on a U-curve.} Putting equal weight on thin midday slices pays
more impact per unit than weighting to the liquid open and close.
\begin{solution}
VWAP weights to where liquidity is---lower impact per slice. It is the right default
when the volume curve is predictable. TWAP is the right default when it is not (quiet
days, news, or deliberately hiding your intent): VWAP's predictability is also a
signature your counterparties can read.
\end{solution}
\end{example}

\begin{figure}[htbp]\centering
\begin{tikzpicture}
\begin{axis}[width=0.86\linewidth,height=0.34\linewidth,
    xlabel={slice},ylabel={\% of day volume},
    xmin=0,xmax=48,ymin=0,
    tick style={draw=none},axis line style={draw=slate},font=\scriptsize]
  \addplot[fill=slate!40,draw=slate] file{figures/data/VOL.tex};
  \node[font=\scriptsize,text=slate] at (axis cs:24,3.4) {U-shape: heavy open \& close};
\end{axis}
\end{tikzpicture}

\caption{The U-shaped intraday volume curve VWAP weights to.}
\label{fig:volume}
\end{figure}

\begin{figure}[htbp]\centering
\begin{tikzpicture}
\begin{axis}[width=0.86\linewidth,height=0.34\linewidth,
    xlabel={slice},ylabel={schedule weight (\%)},
    xmin=0,xmax=48,ymin=0,ymax=7,
    tick style={draw=none},axis line style={draw=slate},font=\scriptsize,legend style={font=\scriptsize,draw=none}]
  \addplot[accent,very thick,domain=0:48] {100/48};
  \addlegendentry{TWAP (equal)}
  \addplot[good,very thick] file{figures/data/VOL.tex};
  \addlegendentry{VWAP (vol-weighted)}
\end{axis}
\end{tikzpicture}

\caption{TWAP (flat) vs VWAP (volume-weighted) schedules over the same 48 slices.}
\label{fig:sched}
\end{figure}

\section{Implementation shortfall}\label{sec:is}

IS measures the trade against the mid at decision time:
$\mathrm{IS}=\mathrm{impact}+\mathrm{timing\ drift}+\mathrm{spread}+\mathrm{fees}$.
\texttt{is\_shortfall.py} on $Q=100{,}000$ (10\% ADV):
\begin{lstlisting}[numbers=none]
instant IS: +63.2bp
split  IS: +44.7bp impact + +7.6bp drift on half = +48.5bp
splitting trades speed for risk: less impact, more timing exposure
\end{lstlisting}

\begin{example}\label{ex:tradeoff}
\textbf{The fundamental trade.} Instant execution pays 63bp of impact; splitting pays
45bp of impact plus 8bp of drift.
\begin{solution}
Speed is a hedge against adverse drift; patience is a hedge against impact. The optimal
schedule sits where marginal impact saved equals marginal timing risk added---and that
point is different for every volatility regime, which is why ``the right tempo'' is a
risk-management question, not a comfort question.
\end{solution}
\end{example}

\begin{figure}[htbp]\centering
\begin{tikzpicture}
\begin{axis}[width=0.86\linewidth,height=0.34\linewidth,
    xlabel={slice},ylabel={mid price},
    xmin=0,xmax=48,
    tick style={draw=none},axis line style={draw=slate},font=\scriptsize]
  \addplot[ink,thick] file{figures/data/PX.tex};
  \addplot[red,dashed] coordinates {(0,100)(48,100)};
  \node[font=\scriptsize,text=red] at (axis cs:40,100.03) {arrival mid};
\end{axis}
\end{tikzpicture}

\caption{The walk the schedule trades through: every slice pays the then-current price,
so drift between decision and fill is uncompensated risk.}
\label{fig:path}
\end{figure}

\begin{figure}[htbp]\centering
\begin{tikzpicture}[font=\scriptsize]
  \draw[->,slate] (0,0)--(8.4,0) node[below left=-1pt]{cost decomposition};
  \draw[fill=accent!25,draw=accent] (0.4,0) rectangle (2.2,2.4) node[midway,align=center]{arrival\\ price\\ risk};
  \draw[fill=amber!40,draw=amber] (2.6,0) rectangle (4.4,1.6) node[midway,align=center]{tempo\\ (impact $\downarrow$\\ drift $\uparrow$)};
  \draw[fill=good!30,draw=good] (4.8,0) rectangle (6.6,3.0) node[midway,align=center]{IS =\\ impact +\\ timing +\\ spread/fees};
  \node[align=left,text=slate] at (9.6,1.2) {instant: 63bp\\ split: 44.7bp +\\ 3.8bp drift};
\end{tikzpicture}

\caption{IS decomposition: tempo moves cost between the impact bucket and the timing-risk
bucket; it does not delete cost.}
\label{fig:is}
\end{figure}

\section{The trader's calculus}\label{sec:trader}

Execution is a risk-management problem disguised as a mechanics problem. Decide the
urgency (how much drift hurts), the size (how much impact costs), and the visibility (who
is watching) before touching the keyboard---the order type follows from those three.

\begin{example}\label{ex:capacity}
\textbf{Capacity is an execution number.} A strategy returns 15\%/yr on 5bp of daily edge;
its ADV footprint is 8\%.
\begin{solution}
The impact at 8\% ADV is $\sigma\sqrt{0.08}\approx57$bp/day---more than ten times the edge.
The strategy is not a \$100M strategy with execution overhead; it is a \$8M strategy at
best, and capacity should be quoted with the same confidence interval as the alpha.
\end{solution}
\end{example}

\begin{checkpoint}[Execution audit]
\begin{itemize}
\item What does concavity buy a big order? (Slicing lowers total impact.)
\item When is TWAP preferable to VWAP? (Unpredictable volume / hiding intent.)
\item What are the four pieces of IS? (Impact, timing drift, spread, fees.)
\end{itemize}
\end{checkpoint}

\section{Recap}\label{sec:execrecap}

Market orders buy certainty, limit orders sell it. Impact is roughly square-root in
participation, so size gets sliced and tempo trades impact for drift. The honest cost of
any strategy is the execution-adjusted number, not the backtest number---and capacity is
what's left when the edge exceeds the print.

\begin{table}[htbp]\centering\small
\caption{Chapter toolbox.}
\begin{tabular}{@{}lll@{}}
\toprule
Concept & Rule & Gotcha\\
\midrule
Market vs limit & certainty vs selection & limit fills are adverse\\
Impact & $c\sim\sigma\sqrt{q/V}$ & capacity = edge/impact\\
TWAP/VWAP & flat vs volume-weighted & VWAP is readable\\
IS & impact + drift + spread + fees & tempo re-mixes, doesn't delete\\
\bottomrule
\end{tabular}
\end{table}

\begin{example}[The execution frontier]\label{ex:frontier}
\textbf{Strategy.} You must move 10\% of ADV. Faster trading pays
more impact but finishes sooner; slower pays less but carries
timing risk. Compute both legs across a participation grid under
the chapter's $\sqrt{\cdot}$ model.

\textbf{Solution.} \texttt{part\_curve.py} prints:
\begin{verbatim}
part   2%: impact  29.7bp  duration  5.0d  1-sigma risk 670.8bp
part   5%: impact  47.0bp  duration  2.0d  1-sigma risk 424.3bp
part  10%: impact  66.4bp  duration  1.0d  1-sigma risk 300.0bp
part  20%: impact  93.9bp  duration  0.5d  1-sigma risk 212.1bp
part  40%: impact 132.8bp  duration  0.2d  1-sigma risk 150.0bp
impact falls with slower trading, risk grows -- the execution frontier
\end{verbatim}
Each leg answers a different prayer: slow trading halves the
impact ($132\to30$bp) but doubles the timing risk
($150\to670$bp). There is no free lunch --- only a frontier.

\textbf{Significance.} This is why ``just TWAP it'' is a decision,
not a default: the schedule chooses where on the frontier you sit.
The desk's risk aversion picks the point, and the CIO's mandate
picks the risk aversion --- execution policy is portfolio choice
in miniature.
\end{example}

\section{Common mistakes}
\begin{itemize}
\item \textbf{``Impact is linear.''} It is concave; big prints cost disproportionately.
\item \textbf{VWAP on an anomalous day.} Weighting to predicted volume when today's is not
  U-shaped increases cost.
\item \textbf{Counting impact but not drift.} Slow schedules pay timing risk.
\item \textbf{Backtests at mid.} Filling at mid in simulation is an execution subsidy.
\item \textbf{Ignoring adverse selection on limits.} Passive fills happen when informed
  flow is the other side.
\item \textbf{Capacity confusion.} ``The strategy does 15\%'' is meaningless without the
  footprint in ADV.
\item \textbf{Minimising impact alone.} Slower trading
always lowers impact but multiplies timing risk
(Example~\ref{ex:frontier}) --- optimising one leg maximises the
other.
\item \textbf{Treating TWAP as neutral.} A uniform schedule is
still a bet: you hold market risk for the whole window and pay
the average, not the best, price.
\end{itemize}

\section{Exercises}
\begin{exercise}[easy]
What does a market order pay that a limit order does not?
\end{exercise}
\begin{solution}
Half-spread plus impact for certainty; the limit pays adverse selection instead.
\end{solution}

\begin{exercise}[easy]
Write the square-root impact law and define $V$.
\end{exercise}
\begin{solution}
$c(q)\approx\sigma\sqrt{q/V}$, $V$ = daily volume (ADV).
\end{solution}

\begin{exercise}[easy]
At $\sigma=2\%$ and 4\% of ADV, what is the predicted impact?
\end{exercise}
\begin{solution}
$0.02\times\sqrt{0.04}=40$bp.
\end{solution}

\begin{exercise}[easy]
Name the four components of implementation shortfall.
\end{exercise}
\begin{solution}
Impact, timing drift, spread, fees.
\end{solution}

\begin{exercise}[easy]
Why is TWAP sometimes chosen over VWAP despite higher predicted cost?
\end{exercise}
\begin{solution}
Predictability: VWAP's U-weighting is a signature; TWAP hides intent and is robust when
the day's volume shape is anomalous.
\end{solution}

\begin{exercise}
Show the marginal cost under the square-root law declines with quantity.
\end{exercise}
\begin{solution}
$c'(q)\propto q^{-1/2}$: each additional unit is cheaper than the last in relative terms,
which is why slicing pays.
\end{solution}

\begin{exercise}
A 20bp daily edge trades 12\% of ADV. Is it viable?
\end{exercise}
\begin{solution}
$\sigma\sqrt{0.12}\approx69$bp impact---edge consumed $\sim3\times$; not viable at that
footprint.
\end{solution}

\begin{exercise}
Why does adverse selection make limit fills ``expensive'' even when they are free?
\end{exercise}
\begin{solution}
You are filled when informed traders sell to you---the next move is against you on average.
\end{solution}

\begin{exercise}
Derive the optimal number of slices $n$ minimizing total cost under the square-root law
and linear drift risk.
\end{exercise}
\begin{solution}
Cost $\approx \frac{1}{\sqrt n}\cdot\text{impact} + n\cdot\text{drift-risk}$; minimize to
get $n^\star\sim(\text{impact}/\text{drift})^{2/3}$.
\end{solution}

\begin{exercise}
What is POV (percent-of-volume) and when does it beat a fixed schedule?
\end{exercise}
\begin{solution}
Participation tracks realized volume; beats fixed schedules when volume is unpredictable
but you want constant market share.
\end{solution}

\begin{exercise}
A fill at mid in simulation cost 0bp; the live fill cost 14bp. Name the two error sources.
\end{exercise}
\begin{solution}
No spread/impact modeled (simulation subsidy) and possibly favorable-fill bias in the
simulator's matching.
\end{solution}

\begin{exercise}
How does urgency (alpha half-life) change the optimal tempo?
\end{exercise}
\begin{solution}
Short-lived alpha pushes toward faster schedules---pay impact, avoid drift that eats the
signal.
\end{solution}

\begin{exercise}[hard]
Estimate the capacity of a 10bp/day edge trading $\sigma=2\%$ names where you cap impact
at 30\% of edge.
\end{exercise}
\begin{solution}
Need $\sigma\sqrt{q/V}\le3$bp $\Rightarrow q/V\le(0.003/0.02)^2=2.25\%$ ADV---capacity is
2.25\% of daily volume in position turnover.
\end{solution}

\begin{exercise}[hard]
Explain why the square-root (concave) law makes ``one big print'' dominated by ``$n$
slices'' even though each slice is read by the market.
\end{exercise}
\begin{solution}
$\sum_i c(q_i)$ with $q_i=q/n$ gives total $n\cdot c(q/n)=c(q)/\sqrt n$---the concavity,
not stealth, does the work.
\end{solution}

\begin{exercise}[hard]
Formalize the IS drift term: is it a cost or a variance, and does it change the optimal
schedule vs a cost-only model?
\end{exercise}
\begin{solution}
Drift is a variance with mean zero---a cost-only optimizer slices minimally (zero risk),
a mean-variance one trades mean impact against drift variance; optimal tempo is finite.
\end{solution}

\begin{exercise}[easy]
From the frontier: what is the impact saving from 40\% down to
2\% participation, and what does it cost?
\end{exercise}
\begin{solution}
$132.8-29.7\approx103$bp saved; duration goes $0.2$d $\to 5$d so
1$\sigma$ timing risk grows $150\to671$bp --- the saving is
never free.
\end{solution}
\begin{exercise}[medium]
A PM insists ``trade everything today, alpha decays.'' Which
frontier point does that choose, and is it optimal if alpha is
flat?
\end{exercise}
\begin{solution}
The fast end ($\ge40\%$): optimal only when the decay cost of
waiting exceeds the impact premium --- with flat alpha it is
the worst point on the curve.
\end{solution}
\begin{exercise}[hard]
Formalise the trade-off as impact $+ \lambda\cdot$risk and locate
the optimum participation on the printed grid at $\lambda=1$.
What changes at $\lambda=0.1$ vs $\lambda=10$?
\end{exercise}
\begin{solution}
At $\lambda=1$ the minimum of impact$+$risk sits near 5--10\%
($47+424$ vs $66+300$ --- 10\% wins at $366$bp). $\lambda=0.1$
pushes slower (risk cheap, impact dominates); $\lambda=10$
pushes faster.
\end{solution}

\section*{Notes}
Execution is where research meets the tape. The honest edge is the backtest edge minus
the square root of your footprint; everything else is negotiation between impact and
impatience. The traders who get this right are invisible; the ones who don't fund the
ones who do.

\section*{Further reading}
\begin{itemize}
\item Cartea, Jaimungal \& Penalva, \emph{Algorithmic and High-Frequency Trading} ---
Ch.~6--7 on optimal execution:

\url{https://www.cambridge.org/core/books/algorithmic-and-highfrequency-trading/}
\item Almgren \& Chriss, ``Optimal Execution of Portfolio Transactions'' --- the model:

\url{https://www.math.nyu.edu/~chriss/optimal.htm}
\item L.~de Prado lectures --- microstructure and execution slides:

\url{https://www.quantresearch.org/Lectures.htm}
\item Binance Academy, ``What Is Market Impact?'' --- practitioner's explainer:
\url{https://academy.binance.com/}
\item Kraken order-types documentation --- the mechanics live:

\url{https://support.kraken.com/hc/en-us/sections/200577136-Order-types}
\end{itemize}

\begin{rulebox}[Watch instead]
\begin{itemize}
\item ``Initial Return: LOB'' --- Initial Return (18:56). Order flow mechanics; re-verified
2026-10-11.
\url{https://www.youtube.com/watch?v=BJt6PGKeO9I}
\item ``INFORMS Market Microstructure'' --- INFORMS (1:28:22). The institutional view;
re-verified 2026-10-11.
\url{https://www.youtube.com/watch?v=XoBjQqMmKoM}
\end{itemize}
\end{rulebox}
